%--------------------------------------------------------------------%
% Page          : Abstract (Indonesian)
%--------------------------------------------------------------------%

\clearpage

\addcontentsline{toc}{chapter}{ABSTRAK}

\section*{\large \centering \MakeUppercase 
    {Abstrak}
}

\begin{center}

    \justifying \normalsize

    \textbf{Pengambilan Keputusan dan Pengaturan Skalabilitas Elastis pada Lingkungan Cloud yang Heterogen dengan Berbasiskan pada Prediksi Workload}

    \vspace{0.5cm}

    \textbf{Muhammad Zahid Masruri / NRP 6025222041}

    \vspace{0.5cm}

    Arsitektur \emph{cloud} modern telah mengubah cara aplikasi web dikembangkan, digunakan, dan dikelola, menawarkan fleksibilitas, skalabilitas, dan efisiensi yang signifikan. Evolusi ini mencakup pergeseran dari infrastruktur tradisional seperti mesin virtual ke teknologi kontainer dan model \emph{serverless}. \emph{Kubernetes}, sebagai sistem manajemen kluster berbasis kontainer, populer karena kemampuannya untuk skala berdasarkan permintaan. Namun, algoritma skalabilitas tradisional sering gagal menangani lonjakan lalu lintas mendadak, yang menyebabkan peningkatan latensi dan alokasi sumber daya yang tidak efisien. Sebagai solusi, komputasi \emph{serverless} atau \emph{Function as a Service (FaaS)} menyediakan skalabilitas otomatis dan efisiensi biaya dengan mengenakan biaya hanya berdasarkan penggunaan aktual, meskipun model ini menimbulkan biaya tinggi untuk beban kerja konstan.

    Penelitian ini mengeksplorasi integrasi antara \emph{Kubernetes} dan komputasi \emph{serverless} untuk mengoptimalkan kinerja dan efisiensi biaya aplikasi \emph{cloud}. Meskipun beberapa penelitian telah fokus pada prediksi alokasi sumber daya, masih terdapat tantangan dalam pengelolaan lonjakan lalu lintas dan alokasi sumber daya yang belum sepenuhnya teratasi. Penelitian ini mengusulkan strategi distribusi lalu lintas yang memanfaatkan prediksi beban kerja berbasis GRU dan mekanisme reservasi sumber daya untuk memastikan alokasi yang optimal dan pemenuhan \emph{tail latency}. Dengan menggabungkan keunggulan dari \emph{Kubernetes} dan \emph{serverless}, penelitian ini bertujuan untuk menciptakan arsitektur \emph{cloud} yang lebih tangguh, responsif, dan efisien biaya.

    \vspace{0.5cm}

    \textbf{Kata Kunci:} \emph{Cloud Computing, Gated Recurrent Unit, Serverless, Function as A Service, Container, Kubernetes}

\end{center}

\clearpage
